
Continue the proof by induction on the rank @@M0@@ of @@M1@@. Suppose the result proved for torsors under split tori of rank @@M2@@. Write @@M3@@, with @@M4@@ and @@M5@@; put @@M6@@ and @@M7@@, so that @@M8@@. Let @@M9@@; then @@M10@@ is a torsor under @@M11@@, and @@M12@@ is a torsor over @@M13@@ with group @@M14@@ and algebra the inverse image of @@M15@@ on @@M16@@. By the preceding argument, @@M17@@ is the quotient of @@M18@@ by the subgroup generated by the @@M19@@ for @@M20@@. By the induction hypothesis, @@M21@@ is the quotient of @@M22@@ by the subgroup generated by the @@M23@@ for @@M24@@. Taking account of the relation
@@M25@@
@@M26@@ is therefore the quotient of @@M27@@ by the subgroup generated by the @@M28@@ for @@M29@@, or equivalently for @@M30@@ running through a system of generators of @@M31@@.

\sgaStableTarget{sga6.e09.cor.2-3}\textbf{Corollary 2.3.} \uline{With the notation and hypotheses of} 2.1, \uline{suppose that} @@M32@@ \uline{is regular}; \uline{then one has a canonical isomorphism}:
@@M33@@
\uline{where} @@M34@@ \uline{is the ideal of} @@M35@@ \uline{generated by the} @@M36@@ \uline{for} @@M37@@ \uline{running through a system of generators of} @@M38@@.

If @@M39@@ is regular, the canonical homomorphism @@M40@@ is an isomorphism (IV 2.5). Likewise, @@M41@@ is then regular, so @@M42@@ is an isomorphism. The corollary therefore follows immediately from 2.2.

